%---------------------------------------------------------------
\chapter{\babEmpat}
%---------------------------------------------------------------
%---------------------------------------------------------------
\section{Simpulan}
%---------------------------------------------------------------

Pelaksanaan kerja magang sebagai \textit{Full-Stack Developer} di PT Pratama
Nusantara Sakti selama periode 05 Januari 2026 hingga 20 Mei 2026 telah berhasil
diselesaikan dengan keluaran utama berupa sistem \textit{HR Safety Campus},
yaitu sebuah aplikasi administrasi tenaga kerja kontrak berbasis web yang
dibangun menggunakan kerangka kerja Laravel. Berdasarkan seluruh proses yang
telah dilalui, dapat ditarik beberapa simpulan sebagai berikut:

\begin{enumerate}
    \item Antarmuka sistem registrasi tenaga kerja berbasis web berhasil
    dibangun menggunakan \textit{Tailwind CSS} untuk menggantikan antarmuka
    Microsoft Access yang lama. Antarmuka baru ini bersifat responsif dan ringan
    sehingga dapat diakses dengan baik oleh tim IT di lapangan melalui berbagai
    perangkat.

    \item Arsitektur \textit{Role-Based Access Control} (RBAC) berhasil
    diimplementasikan pada logika \textit{backend} untuk membatasi akses menu
    sesuai peran masing-masing pengguna. Melalui penerapan prinsip
    \textit{least privilege}, sistem berhasil memisahkan hak akses antara
    Administrator dengan pengguna operasional, sehingga celah manipulasi pada
    proses pencetakan surat izin dan klaim transportasi dapat diperkecil secara
    signifikan.

    \item Proses pembersihan data (\textit{data cleansing}) dan migrasi data
    tenaga kerja dari sistem lama ke basis data MySQL berhasil dilaksanakan.
    Tantangan utama berupa duplikasi Nomor Induk Kependudukan (NIK) berhasil
    diatasi melalui skrip seleksi data khusus yang memisahkan rekaman ganda
    serta rekaman yang tidak valid sebelum proses unggah, dengan tetap melakukan
    konfirmasi kepada pihak klien guna memastikan tidak ada data sah yang
    terhapus.

    \item Mekanisme validasi data berhasil diperketat melalui sistem, baik pada
    fitur klaim transportasi yang menerapkan aturan bisnis berlapis maupun pada
    fitur persetujuan karyawan yang mengevaluasi kelengkapan data secara
    otomatis. Mekanisme ini efektif dalam meminimalisasi kesalahan input
    (\textit{human error}) dan menutup celah praktik kecurangan klaim oleh pihak
    vendor.
\end{enumerate}

Secara keseluruhan, keenam fitur utama dalam sistem \textit{HR Safety Campus},
yaitu Registrasi Tenaga Kerja, Klaim Transport, Surat Izin Keluar, Manajemen
Tarif Transport, Persetujuan Karyawan, dan Manajemen Pengguna, telah berhasil
dikembangkan dan diintegrasikan. Sistem ini menjawab kebutuhan perusahaan akan
proses administrasi yang lebih cepat, akurat, dan terpusat dibandingkan sistem
lokal berbasis Microsoft Access yang digunakan sebelumnya.

%---------------------------------------------------------------
\section{Saran}
%---------------------------------------------------------------

Meskipun sistem \textit{HR Safety Campus} telah berhasil dikembangkan dan
memenuhi kebutuhan dasar perusahaan, terdapat sejumlah saran yang dapat menjadi
bahan pertimbangan untuk pengembangan lebih lanjut, antara lain:

\begin{enumerate}

    \item Mekanisme validasi data pada saat \textit{import} massal dapat
    diperkuat dengan menambahkan lapisan validasi otomatis di sisi sistem,
    seperti pemeriksaan format NIK dan konsistensi tanggal, sehingga data
    bermasalah dapat terdeteksi lebih dini sebelum masuk ke basis data utama
    tanpa harus bergantung pada proses pemeriksaan manual.

    \item Penerapan lapisan keamanan tambahan seperti \textit{rate limiting}
    dan pencatatan \textit{audit log} yang lebih komprehensif dapat
    dipertimbangkan untuk memperkuat sistem pertahanan terhadap akses tidak sah,
    mengingat sistem menyimpan data pribadi ribuan tenaga kerja yang bersifat
    sensitif.

    \item Untuk pelaksanaan kerja magang di masa mendatang, disarankan agar
    dokumentasi terhadap sistem lama dilakukan secara lebih menyeluruh pada tahap
    awal. Hal ini akan membantu mempercepat proses pemahaman alur kerja dan
    mengurangi waktu yang dibutuhkan pada tahap observasi dan migrasi data.
\end{enumerate}
